\documentclass[11pt]{article}

\usepackage[a4paper,margin=2.6cm]{geometry}
\usepackage[utf8]{inputenc}
\usepackage[T1]{fontenc}
\usepackage{amsmath,amssymb,amsthm}
\usepackage{graphicx}
\usepackage{booktabs}
\usepackage{caption}
\usepackage{enumitem}
\usepackage{xcolor}
\usepackage[colorlinks=true,linkcolor=blue!50!black,citecolor=blue!50!black,urlcolor=blue!50!black]{hyperref}
\usepackage{mathtools}

\newcommand{\dd}{\mathrm{d}}
\newcommand{\pdv}[2]{\dfrac{\partial #1}{\partial #2}}
\newcommand{\ppdv}[2]{\dfrac{\partial^2 #1}{\partial #2^2}}
\newcommand{\avg}[1]{\langle #1 \rangle}
\newcommand{\kB}{k_{\mathrm{B}}}
\newcommand{\Var}{\operatorname{Var}}

\title{\Large Brownian Dynamics of a Trapped Particle in Optical Tweezers:\\[2pt]
\large From the Harmonic Trap to Non-Conservative Forces, Escape, and Capture}
\author{Joshua Cox\\[2pt]
\small Winter Research Project, University of Queensland\\
\small Supervisor: Alex Stilgoe}
\date{July 2026}

\begin{document}
\maketitle

\begin{abstract}
\noindent A dielectric particle held in optical tweezers and immersed in a viscous fluid is, on the timescales of interest, an overdamped Brownian particle moving in an external force field. This report develops the stochastic and statistical-mechanical theory of such a particle through four stages of increasing physical realism. Phase 1 derives the governing equations from first principles -- the inertial Langevin equation, the overdamped (Smoluchowski) limit, the fluctuation--dissipation theorem, the exact Ornstein--Uhlenbeck solution, and the equivalent Fokker--Planck description -- and validates two independent numerical methods against the closed-form Boltzmann solution. Phase 2 extends the theory to the axisymmetric $(r,z)$ geometry of a real Gaussian-beam trap, where the stationary state factorises exactly into a Rayleigh radial marginal and a Gaussian axial marginal. Phase 3 introduces the non-conservative scattering (radiation-pressure) force; the resulting force field has a nonzero curl, detailed balance breaks, and the stationary state supports a circulating, divergence-free probability current -- a Brownian vortex -- with no closed-form density. Phase 4 replaces the infinitely deep harmonic well with the finite-depth Gaussian-beam potential and develops the closed-form first-passage theory of trap escape (the exact one-dimensional mean-first-passage-time quadrature and its deep-well Kramers/Arrhenius limit) and of particle capture from the bulk (the Debye--Smoluchowski rate). Combining the two yields the central applied result: retention time rises exponentially with beam power while time-to-contamination falls, so the single-particle holding window is maximised at a well-defined compromise power. Throughout, every result is verified two independent ways -- Brownian dynamics against a direct Fokker--Planck solve, both against the analytic theory wherever one exists -- with numerical agreement at or below the percent level; the numerical details are kept to a minimum here in favour of the closed-form theory.
\end{abstract}

\tableofcontents

% ======================================================================
\section{Introduction and physical setting}
% ======================================================================

Optical tweezers confine a small dielectric particle near the focus of a tightly focused laser beam using the gradient force associated with the induced dipole in an inhomogeneous field \cite{JonesMaragoVolpe2015}. Near the focus, to leading order in the displacement from equilibrium, the trapping potential is harmonic,
\begin{equation}
U(\mathbf{x}) = \tfrac{1}{2}k\,x^2 + \dots,
\label{eq:harmonic-potential}
\end{equation}
and the particle is simultaneously subject to two effects of the surrounding fluid that are not independent of one another: viscous drag, which damps its motion, and thermal agitation from molecular collisions, which drives it. The central theoretical fact underlying everything that follows is that these two effects are two manifestations of the same microscopic process, and are therefore linked quantitatively by the \emph{fluctuation--dissipation theorem}.

This report is the write-up of a staged research project (Table~\ref{tab:phases}) whose scientific aim is to use a single trapped nanoparticle as a volume-exploring probe of its local optical and force environment. Because a single trapped particle has no inter-particle interactions, its equilibrium position density is simply the barometric (ideal-gas) density in the trap potential -- but once the optical force field is allowed to have a non-conservative (scattering) component, that density no longer has closed form, and the only viable route to it is to integrate the governing equation numerically. The project's organising methodological principle is therefore that \emph{every result is checked two independent ways}: Brownian dynamics (stochastic trajectory simulation) against a direct numerical solution of the Fokker--Planck equation, both checked in turn against whatever analytic result is available. Phase 1 establishes this three-way check in the simplest non-trivial setting, the one-dimensional conservative harmonic trap, where the Ornstein--Uhlenbeck/Boltzmann solution is exact ground truth. Phase 2 carries it into the axisymmetric $(r,z)$ geometry, still with an exact factorised stationary state. Phase 3 deliberately gives up the analytic leg -- the non-conservative steady state has no closed form -- so the two-method cross-check becomes the \emph{only} route to a trusted answer. Phase 4 then restores the analytic leg from an unexpected direction: escape from and capture into a finite-depth well both have exact closed forms in the overdamped limit (the mean-first-passage-time quadrature and the Debye--Smoluchowski rate), so the finite-depth physics is again a three-way check.

\begin{table}[h]
\centering
\begin{tabular}{@{}cll@{}}
\toprule
Phase & Description & Status \\
\midrule
0 & Setup (parameters, conventions) & done \\
1 & 1D conservative test rig (\S\ref{sec:langevin}--\S\ref{sec:numerics1d}) & done \\
2 & Axisymmetric $(r,z)$ conservative trap (\S\ref{sec:phase2}) & done \\
3 & Non-conservative scattering forces, Brownian vortex (\S\ref{sec:phase3}) & done \\
4 & Finite-depth trap: escape, capture, and the power trade-off (\S\ref{sec:phase4}) & done \\
5 & Write-up (\emph{this report}) & in progress \\
\bottomrule
\end{tabular}
\caption{Project roadmap. This report presents the mathematics and physics of Phases 1--4, with the numerical validation summarised compactly in \S\ref{sec:validation}.}
\label{tab:phases}
\end{table}

% ======================================================================
\section{Stochastic dynamics of the trapped particle}
\label{sec:langevin}
% ======================================================================

\subsection{Hydrodynamic drag}
\label{sec:drag}

A sphere of radius $a$ moving slowly through a fluid of viscosity $\eta$ experiences a linear (Stokes) drag force $F_{\rm drag} = -\gamma v$, with
\begin{equation}
\gamma = 6\pi \eta a,
\label{eq:stokes}
\end{equation}
obtained by solving the incompressible Navier--Stokes equations in the creeping-flow (zero Reynolds number) limit,
\begin{equation}
\mathrm{Re} = \frac{\rho U a}{\eta} \ll 1,
\end{equation}
subject to a no-slip boundary condition at the sphere's surface and quiescent flow at infinity. In this limit the inertial (nonlinear advective) term in Navier--Stokes is negligible relative to the viscous term, leaving the linear Stokes equations $\eta \nabla^2 \mathbf{u} = \nabla p$, $\nabla\cdot\mathbf{u}=0$, whose solution around a translating sphere yields Eq.~\eqref{eq:stokes}. This is valid provided the particle is large compared to the solvent's molecular scale (so continuum hydrodynamics applies) but small/slow enough that fluid inertia is negligible -- both hold comfortably for a sub-micron bead in water at the speeds considered here.

\subsection{The inertial Langevin equation and fluctuation--dissipation}

The full (inertial) Langevin equation for a bead of mass $m$ in the harmonic trap of Eq.~\eqref{eq:harmonic-potential}, immersed in a viscous fluid at temperature $T$, is
\begin{equation}
m\,\dot v = -\gamma v - kx + \sqrt{2\gamma \kB T}\,\xi(t),
\label{eq:full-langevin}
\end{equation}
where $\xi(t)$ is a delta-correlated Gaussian white-noise process, $\avg{\xi(t)\xi(t')} = \delta(t-t')$, $\avg{\xi(t)}=0$. The noise amplitude $\sqrt{2\gamma \kB T}$ is not a free parameter: it is fixed relative to the drag coefficient $\gamma$ by the fluctuation--dissipation theorem, which states that the strength of a system's response to a small perturbation (here, the drag coefficient governing the decay of an impulse) is quantitatively tied to the strength of its spontaneous equilibrium fluctuations (here, the thermal noise driving diffusion). Physically, both $\gamma$ and the noise originate in the same random collisions between the bead and solvent molecules: the theorem is what guarantees that Eq.~\eqref{eq:full-langevin}, evolved for long times with no external drive, relaxes to the correct Boltzmann equilibrium at temperature $T$ rather than to some other, unphysical steady state.

\subsection{The overdamped (Smoluchowski) limit}

The momentum relaxation time of Eq.~\eqref{eq:full-langevin} is $m/\gamma$. For the default parameters of this project (a $0.5\,\mu\mathrm{m}$-radius silica bead in water, Table~\ref{tab:params}), $m/\gamma \sim 10^{-7}\,\mathrm{s}$, many orders of magnitude shorter than the position relaxation time $\tau = \gamma/k \sim 10^{-2}\,\mathrm{s}$ set by the trap. On the timescales of interest, the velocity has therefore already equilibrated to the local force balance; formally, one takes $m/\gamma \to 0$ and drops the inertial term $m\dot v$ entirely from Eq.~\eqref{eq:full-langevin} (the \emph{overdamped} or \emph{Smoluchowski} limit). What remains is a first-order stochastic differential equation (SDE) directly for position,
\begin{equation}
0 = -\gamma \dot x - kx + \sqrt{2\gamma \kB T}\,\xi(t)
\quad\Longrightarrow\quad
\dot x = -\frac{k}{\gamma}x + \sqrt{2D}\,\xi(t),
\label{eq:overdamped-sde}
\end{equation}
where the Einstein relation
\begin{equation}
D = \frac{\kB T}{\gamma}
\label{eq:einstein}
\end{equation}
has been used to trade the noise amplitude $\sqrt{2\gamma \kB T}$ for the diffusion coefficient $D$. Equation~\eqref{eq:overdamped-sde} is a linear SDE with additive (state-independent) noise -- an \emph{Ornstein--Uhlenbeck (OU) process}. Two derived scales organise the rest of the analysis:
\begin{align}
\tau &= \frac{\gamma}{k} &&\text{(relaxation time, from force balance } \gamma\dot x = -kx \text{),}
\label{eq:tau}\\
\sigma_x &= \sqrt{\frac{\kB T}{k}} &&\text{(stationary standard deviation, from equipartition).}
\label{eq:sigmax}
\end{align}
Equation~\eqref{eq:sigmax} follows from the equipartition theorem applied to the harmonic well: $\avg{U} = \avg{\tfrac12 k x^2} = \tfrac12 \kB T$ at thermal equilibrium, so $\Var(x) = \kB T/k \equiv \sigma_x^2$.

\subsection{Exact solution of the Ornstein--Uhlenbeck process}
\label{sec:ou-solution}

Equation~\eqref{eq:overdamped-sde} is linear, and so can be solved in closed form with an integrating factor exactly as one solves the deterministic ODE $\dot x = -x/\tau$. Multiplying through by $e^{t/\tau}$ and recognising the left side as a total derivative,
\begin{equation}
\dd\!\left[x(t)\,e^{t/\tau}\right] = \sqrt{2D}\,e^{t/\tau}\,\dd W(t),
\end{equation}
then integrating from $0$ to $t$ with deterministic initial condition $x(0)=x_0$,
\begin{equation}
x(t) = x_0\,e^{-t/\tau} + \sqrt{2D}\int_0^t e^{-(t-s)/\tau}\,\dd W(s).
\label{eq:ou-solution}
\end{equation}
The first term is deterministic and gives the ensemble mean directly, since the stochastic integral (an It\^o integral against a martingale) has zero mean:
\begin{equation}
\avg{x(t)} = x_0\,e^{-t/\tau}.
\label{eq:ou-mean}
\end{equation}
The variance of the stochastic term follows from the It\^o isometry, which converts the variance of a stochastic integral into an ordinary (deterministic) integral of the squared integrand,
\begin{equation}
\Var\!\left[\int_0^t e^{-(t-s)/\tau}\dd W(s)\right] = \int_0^t e^{-2(t-s)/\tau}\,\dd s,
\end{equation}
so that
\begin{equation}
\Var[x(t)] = 2D\int_0^t e^{-2(t-s)/\tau}\,\dd s
= D\tau\left(1-e^{-2t/\tau}\right)
= \sigma_x^2\left(1-e^{-2t/\tau}\right),
\label{eq:ou-variance}
\end{equation}
using $D\tau = (\kB T/\gamma)(\gamma/k) = \kB T/k = \sigma_x^2$ -- i.e., the Einstein relation Eq.~\eqref{eq:einstein} combined with the definition of $\tau$, Eq.~\eqref{eq:tau}. As $t\to\infty$, Eq.~\eqref{eq:ou-variance} correctly limits to the equipartition result $\sigma_x^2$. Because the SDE is linear with a deterministic initial condition and additive Gaussian noise, $x(t)$ is exactly Gaussian at every $t$, not merely asymptotically; the pair (mean, variance) above therefore fully characterises the transition density $P(x,t\mid x_0)$ at all times.

% ======================================================================
\section{The Fokker--Planck (Smoluchowski) description}
\label{sec:fp}
% ======================================================================

\subsection{From the SDE to the forward Kolmogorov equation}

Trajectory-level (Langevin) and density-level (Fokker--Planck) descriptions are two faces of the same stochastic process. For the general overdamped SDE $\dd x = a(x)\,\dd t + b\,\dd W$ with drift $a(x) = -(k/\gamma)x$ and constant diffusion amplitude $b=\sqrt{2D}$, the probability density $P(x,t)$ of an ensemble of trajectories obeys the forward Kolmogorov (Fokker--Planck, or in this overdamped context, \emph{Smoluchowski}) equation
\begin{equation}
\pdv{P}{t} = -\pdv{}{x}\!\left[a(x)P\right] + \frac{1}{2}\ppdv{}{x}\!\left[b^2 P\right]
= \pdv{}{x}\!\left[\frac{k}{\gamma}xP\right] + D\ppdv{P}{x}
\equiv -\pdv{J}{x},
\label{eq:fp}
\end{equation}
which can be obtained rigorously as the leading (second) order truncation of the Kramers--Moyal expansion of the master equation for a continuous Markov process, exact here because the noise is Gaussian and all cumulants beyond the second vanish identically \cite{Gardiner2009}. The probability current
\begin{equation}
J(x,t) = -\frac{k}{\gamma}xP(x,t) - D\,\pdv{P}{x}(x,t)
\label{eq:current}
\end{equation}
decomposes into an advective piece (drift toward the trap centre carries probability with it) and a Fickian diffusive piece (probability spreads down its own gradient). Equation~\eqref{eq:fp}, $\partial_t P = -\partial_x J$, is a continuity equation, formally identical to conservation of mass or charge; integrating over all $x$ shows $\dd/\dd t \int P\,\dd x = 0$ provided $J\to 0$ at the domain boundaries (or at infinity), so that total probability is conserved for all time.

\subsection{Detailed balance and the Boltzmann stationary state}
\label{sec:detailed-balance}

At a stationary state, $\partial_t P = 0$, which by Eq.~\eqref{eq:fp} requires only that $J$ be independent of $x$; for a bounded (or symmetric, decaying) domain this is only consistent with $J \equiv 0$. For a \emph{conservative} force field this zero-current condition is achieved pointwise -- \emph{detailed balance} -- not merely on average: setting $J=0$ in Eq.~\eqref{eq:current} gives the first-order ODE
\begin{equation}
-\frac{k}{\gamma}xP_{\rm eq}(x) = D\,\frac{\dd P_{\rm eq}}{\dd x},
\end{equation}
whose solution is
\begin{equation}
P_{\rm eq}(x) = \frac{1}{Z}\exp\!\left(-\frac{kx^2}{2\kB T}\right), \qquad Z = \int e^{-U(x)/\kB T}\,\dd x = \sqrt{\frac{2\pi \kB T}{k}},
\label{eq:boltzmann}
\end{equation}
i.e. exactly the Boltzmann distribution for $U(x)=\tfrac12 kx^2$, with variance $\kB T/k = \sigma_x^2$. This must agree with the $t\to\infty$ limit of Eq.~\eqref{eq:ou-variance}, and does: the two derivations of $\sigma_x^2$ -- one from the long-time limit of the driven SDE, one from a static equipartition/detailed-balance argument with no dynamics involved -- coincide only because $D\tau = \kB T/k$ exactly, itself a direct consequence of the fluctuation--dissipation relation Eq.~\eqref{eq:einstein}. Their agreement is therefore a non-trivial self-consistency check on the physical model, not just an algebraic coincidence.

This zero-current equilibrium is special to conservative (gradient) force fields. It is exactly the property that fails once a non-conservative force is introduced in \S\ref{sec:phase3}: there, $J\neq 0$ even in the stationary state, giving rise to a circulating \emph{Brownian vortex} probability current with no closed-form density.

% ======================================================================
\section{Numerical methods for the one-dimensional trap}
\label{sec:numerics1d}
% ======================================================================

Two independent numerical solutions of the same physics are implemented and cross-validated: a stochastic (Brownian-dynamics) trajectory simulation of Eq.~\eqref{eq:overdamped-sde}, and a deterministic finite-volume solution of the density equation, Eq.~\eqref{eq:fp}. The machinery built here -- and stress-tested where an exact answer exists -- is reused, direction by direction, in every later phase.

\subsection{Brownian dynamics: Euler--Maruyama}

Euler--Maruyama is the stochastic analogue of the forward-Euler method: truncating the It\^o--Taylor expansion of $\dd x = a(x)\,\dd t + b(x)\,\dd W$ at first order in $\dd t$ gives the update
\begin{equation}
x_{n+1} = x_n - \frac{k}{\gamma}x_n\,\Delta t + \sqrt{2D\Delta t}\,\xi_n, \qquad \xi_n \sim \mathcal{N}(0,1).
\label{eq:em-update}
\end{equation}
In general, Euler--Maruyama is only strong-order $0.5$ accurate, because a state-dependent diffusion amplitude $b(x)$ requires the additional Milstein correction to reach strong order 1. Here $b$ is constant, so the Milstein correction vanishes identically and the scheme is already strong-order 1. The scheme is only conditionally stable: the noise-free part of the update is the linear map $x \mapsto (1-\Delta t/\tau)x$, bounded only for $\Delta t < 2\tau$; this project fixes $\Delta t = \tau/200$. An ensemble of $N$ independent sample paths is a Monte Carlo estimator of the transition density $P(x,t)$, with statistical error shrinking as $N^{-1/2}$; $N=20{,}000$ trajectories are used throughout.

\subsection{Fokker--Planck: the Chang--Cooper spatial scheme}
\label{sec:chang-cooper}

Naive central differencing of the flux \eqref{eq:current} is stable and positivity-preserving only while the local cell P\'eclet number $\mathrm{Pe} = v\,\Delta x/D$ (with $v(x) = -kx/\gamma$ the drift velocity) stays below $2$ -- a condition eventually violated away from the trap centre, where central differencing produces spurious oscillations and negative densities. The Chang--Cooper scheme \cite{ChangCooper1970} avoids this by discretising the flux at each cell face using the \emph{exact local steady-state solution} of Eq.~\eqref{eq:current}: treating the drift as locally constant between adjacent cell centres, the constant-$J$ ODE $vP - D\,\dd P/\dd x = J$ is solved exactly across the face, yielding
\begin{equation}
J_{i+1/2} = \frac{D}{\Delta x}\Big[B(-\mathrm{Pe})\,P_i - B(\mathrm{Pe})\,P_{i+1}\Big], \qquad
B(z) = \frac{z}{e^z - 1},
\label{eq:cc-flux}
\end{equation}
with $B(0)=1$. The Bernoulli function $B$ is strictly positive for every real $z$, which makes the assembled tridiagonal generator $L$ an M-matrix and the scheme positivity-preserving at \emph{any} P\'eclet number: as $|\mathrm{Pe}|\to 0$ it reduces to central differencing, and as $|\mathrm{Pe}|\to\infty$ it degrades gracefully to upwinding. (Equation~\eqref{eq:cc-flux} is algebraically the Scharfetter--Gummel scheme of semiconductor device simulation.) This exponential-fitting property matters later: the escape problem of \S\ref{sec:phase4} lives precisely in the high-P\'eclet region near a steep potential wall. No-flux (reflecting) boundaries are imposed by zeroing the outermost face fluxes, which makes the finite-volume update telescope exactly: total probability is conserved to machine precision by construction.

\subsection{Time integration: Crank--Nicolson}

With $L$ built once, $\dot{\mathbf{P}}=L\mathbf{P}$ is a linear ODE system, advanced with the Crank--Nicolson (trapezoidal) scheme
\begin{equation}
\left(I - \frac{\Delta t}{2}L\right)\mathbf{P}^{n+1} = \left(I + \frac{\Delta t}{2}L\right)\mathbf{P}^n,
\label{eq:cn}
\end{equation}
which is second-order accurate in $\Delta t$ and unconditionally stable for this operator, since $L$'s eigenvalues are real and non-positive and the amplification factor $|(1+\tfrac12\Delta t\lambda)/(1-\tfrac12\Delta t\lambda)| \le 1$ for every $\lambda \le 0$. Each step is one tridiagonal solve (Thomas algorithm, $O(N)$); diagonal dominance -- again the M-matrix property -- guarantees the elimination never encounters a vanishing pivot.

\subsection{Phase-1 validation in brief}

With the default parameters of Table~\ref{tab:params}, both methods were run from $x_0 = 3\sigma_x$ for $50\tau$ and compared against Eqs.~\eqref{eq:ou-mean}, \eqref{eq:ou-variance} and \eqref{eq:boltzmann}. The Brownian-dynamics stationary standard deviation matches the analytic value to $0.28\%$ and the Fokker--Planck one to numerical precision; probability is conserved to $\sim 10^{-12}$ and the density stays non-negative throughout (Figure~\ref{fig:phase1}, and \S\ref{sec:validation}). The two numerical methods -- one stochastic and trajectory-based, one deterministic and density-based, sharing no code path -- thus agree with the exact theory and with each other, which is the licence to trust them in the later regimes where no exact theory exists.

\begin{table}[h]
\centering
\begin{tabular}{@{}llll@{}}
\toprule
Quantity & Symbol & Definition & Value \\
\midrule
Radius & $a$ & --- & $0.5\,\mu\mathrm{m}$ \\
Viscosity & $\eta$ & --- & $10^{-3}\,\mathrm{Pa\,s}$ (water) \\
Temperature & $T$ & --- & $300\,\mathrm{K}$ \\
Trap stiffness & $k$ & --- & $10^{-6}\,\mathrm{N/m}$ ($\sim 1\,\mathrm{pN}/\mu\mathrm{m}$) \\
\midrule
Drag coefficient & $\gamma$ & $6\pi\eta a$ & $9.4248\times10^{-9}\,\mathrm{kg/s}$ \\
Diffusion coefficient & $D$ & $\kB T/\gamma$ & $4.3947\times10^{-13}\,\mathrm{m^2/s}$ \\
Relaxation time & $\tau$ & $\gamma/k$ & $9.42\,\mathrm{ms}$ \\
Stationary std.\ dev. & $\sigma_x$ & $\sqrt{\kB T/k}$ & $64.4\,\mathrm{nm}$ \\
\bottomrule
\end{tabular}
\caption{Baseline physical parameters and derived quantities (\texttt{src/params.py}): a silica microsphere in water at room temperature, the regime of Volpe \& Volpe's tutorial treatment \cite{VolpeVolpe2013}. Phases 2--3 reuse this regime with an anisotropic stiffness (Table~\ref{tab:params2}); Phase 4 re-parametrises the trap by beam power for a nanodiamond (\S\ref{sec:power}).}
\label{tab:params}
\end{table}

\begin{figure}[h]
\centering
\includegraphics[width=0.92\textwidth]{figures/phase1_validation.png}
\caption{Phase-1 validation (\texttt{scripts/validate\_phase1.py}): Brownian-dynamics and Fokker--Planck relaxation of the mean and variance against the analytic OU curves \eqref{eq:ou-mean}--\eqref{eq:ou-variance}, and both stationary densities against the Boltzmann distribution \eqref{eq:boltzmann}.}
\label{fig:phase1}
\end{figure}

% ======================================================================
\section{Phase 2: the axisymmetric $(r,z)$ trap}
\label{sec:phase2}
% ======================================================================

\subsection{Anisotropic harmonic potential and separability}

A real Gaussian-beam focus is stiffer transverse to the beam axis than along it: the beam waist $w_0$ sets the steep transverse intensity gradient, while the much longer Rayleigh range $z_R$ sets the weaker axial one (this geometrical statement is made quantitative in \S\ref{sec:gaussian-beam}, where both stiffnesses are derived from the beam parameters). The harmonic approximation of the trap therefore carries two independent spring constants,
\begin{equation}
U(r,z) = \tfrac12 k_r r^2 + \tfrac12 k_z z^2, \qquad k_z < k_r,
\label{eq:U-axisym}
\end{equation}
with $r^2 = x^2 + y^2$. This project takes $k_r$ equal to the Phase-1 $k$ and $k_z = k_r/5$, producing the characteristic cigar-shaped (prolate) thermal cloud. The essential structural fact is that Eq.~\eqref{eq:U-axisym} is \emph{separable in Cartesian coordinates}: $U = \tfrac12 k_r x^2 + \tfrac12 k_r y^2 + \tfrac12 k_z z^2$, and since the thermal noise acts independently and isotropically on each coordinate (the same $D$ in all directions -- drag depends on the particle, not the trap), the overdamped SDE decouples into three independent copies of the Phase-1 OU process,
\begin{equation}
\dot x = -\frac{x}{\tau_r} + \sqrt{2D}\,\xi_x, \qquad
\dot y = -\frac{y}{\tau_r} + \sqrt{2D}\,\xi_y, \qquad
\dot z = -\frac{z}{\tau_z} + \sqrt{2D}\,\xi_z,
\label{eq:axisym-sde}
\end{equation}
with $\tau_r = \gamma/k_r$, $\tau_z = \gamma/k_z$, and $\xi_x,\xi_y,\xi_z$ independent white noises. Every Phase-1 result -- the exact Gaussian transition density, the exponential relaxation of mean and variance -- applies verbatim to each coordinate, with its own $(\tau,\sigma)$ pair:
\begin{equation}
\sigma_r = \sqrt{\kB T/k_r}, \qquad \sigma_z = \sqrt{\kB T/k_z}.
\end{equation}

\begin{table}[h]
\centering
\begin{tabular}{@{}lll@{}}
\toprule
Quantity & Definition & Value \\
\midrule
Radial stiffness & $k_r$ & $10^{-6}\,\mathrm{N/m}$ \\
Axial stiffness & $k_z = k_r/5$ & $2\times10^{-7}\,\mathrm{N/m}$ \\
Radial relaxation time & $\tau_r = \gamma/k_r$ & $9.4\,\mathrm{ms}$ \\
Axial relaxation time & $\tau_z = \gamma/k_z$ & $47.1\,\mathrm{ms}$ \\
Per-axis transverse std.\ dev. & $\sigma_r = \sqrt{\kB T/k_r}$ & $64\,\mathrm{nm}$ \\
Axial std.\ dev. & $\sigma_z = \sqrt{\kB T/k_z}$ & $144\,\mathrm{nm}$ \\
\bottomrule
\end{tabular}
\caption{Phase-2 anisotropic trap parameters (\texttt{AxisymTrapParams}); particle and fluid as in Table~\ref{tab:params}.}
\label{tab:params2}
\end{table}

\subsection{The cylindrical Smoluchowski equation}

The Fokker--Planck description does \emph{not} simply inherit the Cartesian decoupling, because the natural coordinates of an axisymmetric problem carry a non-trivial measure. For an axisymmetric density the azimuthal angle is trivial, and the continuity equation for $P(r,z,t)$ in cylindrical coordinates reads
\begin{equation}
\pdv{P}{t} = -\frac{1}{r}\pdv{(r J_r)}{r} - \pdv{J_z}{z},
\qquad
J_r = \frac{F_r}{\gamma}P - D\pdv{P}{r}, \quad
J_z = \frac{F_z}{\gamma}P - D\pdv{P}{z},
\label{eq:fp-axisym}
\end{equation}
with $F_r = -k_r r$ and $F_z = -k_z z$ in the conservative case. The $1/r\,\partial_r(r\,\cdot)$ structure is the cylindrical divergence: probability in the annulus $[r, r+\dd r]$ occupies a volume $\propto r\,\dd r\,\dd z$, so a flux crossing a face at radius $r$ is weighted by the circumference of that face. This Jacobian is precisely what distinguishes the radial problem from two more copies of the 1D equation, and it must be respected both in the analytic marginals below and in the finite-volume discretisation.

\subsection{Stationary state: Rayleigh and Gaussian marginals}
\label{sec:rayleigh}

The force is conservative, so detailed balance holds and the stationary state is the Boltzmann density for Eq.~\eqref{eq:U-axisym}, which factorises exactly because $U$ is separable:
\begin{equation}
P_{\rm eq}(x,y,z) = \frac{1}{Z}\,
e^{-x^2/2\sigma_r^2}\,e^{-y^2/2\sigma_r^2}\,e^{-z^2/2\sigma_z^2}.
\label{eq:boltzmann-axisym}
\end{equation}
The two marginals used for validation are \emph{not} both Gaussian. The axial marginal is Gaussian, exactly as in 1D, since $z$ is a Cartesian coordinate:
\begin{equation}
P(z) = \frac{1}{\sqrt{2\pi}\,\sigma_z}\,e^{-z^2/2\sigma_z^2}.
\end{equation}
The radial marginal is the density of the modulus $r = \sqrt{x^2+y^2}$ of two i.i.d.\ Gaussians. Integrating Eq.~\eqref{eq:boltzmann-axisym} over the angle picks up the $r\,\dd r$ measure,
\begin{equation}
P(r)\,\dd r = \frac{e^{-r^2/2\sigma_r^2}}{2\pi\sigma_r^2}\,\cdot 2\pi r\,\dd r
\quad\Longrightarrow\quad
P(r) = \frac{r}{\sigma_r^2}\,e^{-r^2/2\sigma_r^2},
\qquad r \ge 0,
\label{eq:rayleigh}
\end{equation}
a \emph{Rayleigh} distribution: the same Jacobian that produced the $1/r\,\partial_r(r\,\cdot)$ divergence in Eq.~\eqref{eq:fp-axisym} converts the flat Gaussian into a density that vanishes linearly at the axis and peaks at $r = \sigma_r$. Its second moment is
\begin{equation}
\avg{r^2} = \avg{x^2} + \avg{y^2} = 2\sigma_r^2,
\end{equation}
twice the per-axis variance -- an easy factor to lose, and the first quantity every Phase-2 check is run against. The Rayleigh marginal \eqref{eq:rayleigh} plays an unexpectedly durable role: it survives \emph{unchanged} into the non-conservative Phase 3 (\S\ref{sec:rayleigh-invariant}).

\subsection{Numerical treatment in brief}

The Brownian-dynamics side integrates the three Cartesian SDEs \eqref{eq:axisym-sde} directly, forming $r=\sqrt{x^2+y^2}$ only when reporting -- deliberately sharing no code path (and no coordinate system) with the Fokker--Planck solver it cross-checks. The Fokker--Planck side discretises Eq.~\eqref{eq:fp-axisym} in finite-volume form with Chang--Cooper face weighting in each direction; two points are specific to the geometry. First, the radial grid is \emph{staggered} -- cell centres at $(i+\tfrac12)\Delta r$ -- so the innermost face lands exactly at $r=0$, where its flux weight (the face radius) vanishes identically: the $r=0$ regularity condition falls out of the geometry, with no explicit boundary rule. Second, time-stepping uses Peaceman--Rachford ADI \cite{PeacemanRachford1955}: each step splits into two half-steps, implicit in $r$ then implicit in $z$,
\begin{equation}
\left(I - \tfrac{\Delta t}{2}L_r\right)\rho^{*} = \left(I + \tfrac{\Delta t}{2}L_z\right)\rho^{n},
\qquad
\left(I - \tfrac{\Delta t}{2}L_z\right)\rho^{n+1} = \left(I + \tfrac{\Delta t}{2}L_r\right)\rho^{*},
\label{eq:adi}
\end{equation}
each half-step being a batch of independent 1D tridiagonal solves -- the Phase-1 machinery applied line by line. The splitting is second-order in time and unconditionally stable, since $L_r$ and $L_z$ are individually M-matrix generators.

% ======================================================================
\section{Phase 3: non-conservative forces and the Brownian vortex}
\label{sec:phase3}
% ======================================================================

\subsection{The scattering force and its curl}

A real optical trap is not a pure gradient force. Alongside the gradient (dipole) force, the beam exerts a scattering (radiation-pressure) force: photons scattered forward by the particle transfer momentum \emph{along the beam axis}, with a magnitude proportional to the local beam \emph{intensity} -- which varies \emph{transversely}, tracking the beam profile \cite{JonesMaragoVolpe2015}. Phase 3 models this with the simplest field that captures the mismatch:
\begin{equation}
F_r(r) = -k_r r, \qquad
F_z(r,z) = -k_z z + f_{\rm sc}(r), \qquad
f_{\rm sc}(r) = F_0\, e^{-2r^2/w_0^2},
\label{eq:nc-force}
\end{equation}
a Gaussian-beam-shaped axial push, peaked on-axis, with amplitude $F_0$ and transverse $1/e^2$ width $w_0$. Setting $F_0=0$ recovers Phase 2 exactly, which serves as a built-in regression check.

The combined field is non-conservative. In the axisymmetric meridian plane the only possible curl component is azimuthal, and
\begin{equation}
(\nabla\times\mathbf{F})_\theta
= \pdv{F_r}{z} - \pdv{F_z}{r}
= -f_{\rm sc}'(r)
= \frac{4F_0\,r}{w_0^2}\,e^{-2r^2/w_0^2},
\label{eq:curl}
\end{equation}
since $F_r$ has no $z$-dependence and the only $r$-dependence of $F_z$ is through $f_{\rm sc}$. Equation~\eqref{eq:curl} vanishes identically for $F_0=0$ (and on the axis, where $f_{\rm sc}$ is stationary), and is strictly positive for $r>0$ otherwise. A nonzero curl means \emph{no scalar potential exists} whose gradient reproduces $\mathbf F$: the line integral of $\mathbf F$ around a closed loop in the $(r,z)$ plane is nonzero (by Stokes' theorem it equals the flux of the curl through the loop), so the force does net work on any particle carried around such a loop. That single geometric fact demolishes the entire equilibrium framework of \S\ref{sec:detailed-balance}, as follows.

\subsection{Breakdown of detailed balance: the non-equilibrium steady state}

The stationarity condition for the Fokker--Planck equation \eqref{eq:fp-axisym} is only
\begin{equation}
\nabla\cdot\mathbf{J} = 0,
\label{eq:divJ}
\end{equation}
which the Boltzmann-like ansatz satisfies in the strong, pointwise form $\mathbf J \equiv 0$ (detailed balance) whenever the drift is a gradient. Writing the zero-current condition $\mathbf J = (\mathbf F/\gamma)P - D\nabla P = 0$ as $\nabla \ln P = \mathbf F/(\gamma D) = \mathbf F/\kB T$ makes the obstruction explicit: the left side is a gradient, so the equation is solvable only if $\mathbf F$ is one too. With the curl \eqref{eq:curl} nonzero, no density -- of any functional form -- can zero the current everywhere. The stationary state that the dynamics nevertheless relaxes to is therefore a \emph{non-equilibrium steady state}: the density $\rho_{\rm ss}(r,z)$ is time-independent and satisfies Eq.~\eqref{eq:divJ}, but the current $\mathbf J$ is generically nonzero -- and, being divergence-free in a closed domain with no-flux walls, it must \emph{circulate}. This closed loop of steady probability flux is the \emph{Brownian vortex} \cite{Roichman2008,Sun2009}: the trapped particle endlessly cycles, in the ensemble-averaged sense, around a toroidal roll in the meridian plane, driven by the work the non-conservative force does around each circuit and dissipated into the fluid. In the language of stochastic thermodynamics, the steady state carries a constant housekeeping heat flow; equilibrium is the special case in which that flow -- and with it $\mathbf J$ -- vanishes identically.

Two consequences organise the phase. First, there is no closed-form $\rho_{\rm ss}$: the only route to it is numerical, and the only way to \emph{trust} it is the two-method (Brownian-dynamics vs.\ Fokker--Planck) agreement that the earlier phases validated against exact theory. Second, the current itself becomes the object of interest: the qualitative signature that distinguishes this steady state from an equilibrium one is not the density (which remains unimodal and visually Gaussian-like) but the circulation of $\mathbf J$, with the divergence condition \eqref{eq:divJ} replacing the closed-form density as the quantitative stationarity diagnostic.

\subsection{An exact invariant: the radial marginal stays Rayleigh}
\label{sec:rayleigh-invariant}

One exact analytic statement survives the loss of detailed balance, and it is worth deriving because it provides a closed-form check \emph{inside} the non-conservative regime. In the trajectory picture, the push $f_{\rm sc}(r)$ enters only the $z$-equation of \eqref{eq:axisym-sde}:
\begin{equation}
\dot z = -\frac{z}{\tau_z} + \frac{f_{\rm sc}(r)}{\gamma} + \sqrt{2D}\,\xi_z,
\end{equation}
while the $x$ and $y$ equations are untouched -- they do not involve $z$ or $F_0$ at all. The transverse dynamics $(x,y)$ is therefore an \emph{autonomous} two-dimensional OU process for every value of $F_0$: its marginal law cannot know the push exists. Hence the stationary radial marginal is exactly the Rayleigh density \eqref{eq:rayleigh}, and $\avg{r^2} = 2\sigma_r^2$, for \emph{any} $F_0$ -- even though the joint density $\rho_{\rm ss}(r,z)$ no longer factorises and has no closed form. (The coupling is strictly one-way: $z$ listens to $r$ through $f_{\rm sc}(r)$, but $r$ never listens back -- which is also precisely why the curl \eqref{eq:curl} is nonzero.) This invariant is checked directly at $F_0>0$, a sharper test than only recovering it in the $F_0\to0$ limit.

\subsection{Numerical treatment and the current diagnostic in brief}

On the Brownian-dynamics side the push adds $f_{\rm sc}(r)/\gamma$ to each particle's axial drift, evaluated at its own instantaneous $r$. On the Fokker--Planck side the axial Chang--Cooper coefficients, previously one shared tridiagonal for all radii, become an $r$-dependent field -- one distinct tridiagonal system per radial row in the ADI $z$-half-step, solved with a batched Thomas elimination; the radial operator is unchanged since $F_r$ stays conservative. The steady current $\mathbf J$ is then extracted from the \emph{same} Chang--Cooper face weights the solver's operators are built from -- so the diagnostic measures the flux the solver actually computed, not an independently discretised approximation -- and the stationarity check is the discrete form of Eq.~\eqref{eq:divJ}: $\nabla\cdot\mathbf J \approx 0$ cell by cell, while $\mathbf J$ itself is verifiably nonzero and circulating (Figure~\ref{fig:phase3}).

\begin{figure}[h]
\centering
\includegraphics[width=0.8\textwidth]{figures/phase3_validation.png}
\caption{The Brownian vortex (\texttt{scripts/validate\_phase3.py}): streamlines of the stationary probability current $\mathbf J(r,z)$, coloured by $\log_{10}|\mathbf J|$, for $F_0>0$. The current is nonzero and circulating yet divergence-free -- the defining signature of the non-equilibrium steady state of \S\ref{sec:phase3}.}
\label{fig:phase3}
\end{figure}

% ======================================================================
\section{Phase 4: the finite-depth trap -- escape, capture, and the power trade-off}
\label{sec:phase4}
% ======================================================================

Phases 1--3 model a particle that is trapped \emph{forever}: a harmonic well is infinitely deep, so the escape probability is exactly zero and there is no barrier for an outside particle to cross. The applications Phase 4 targets live in the opposite regime. At the nanoparticle end of optical trapping, particles are hard to see and additional particles can invisibly enter the trap, corrupting measurements -- acutely so for optically trapped nanodiamonds in quantum sensing. Raising the trap power suppresses Brownian escape of the held particle but accelerates the capture of contaminants; lowering it does the reverse. Choosing the operating power is therefore a quantitative trade-off between two first-passage processes, and both, remarkably, have closed-form theory in the overdamped limit.

\subsection{The Gaussian-beam potential and its harmonic limit}
\label{sec:gaussian-beam}

The gradient potential of a focused paraxial Gaussian beam of waist $w_0$ and Rayleigh range $z_R$ is proportional to the local intensity, giving
\begin{equation}
U(r,z) = -\frac{U_0}{s(z)}\,\exp\!\left(-\frac{2r^2}{w_0^2\,s(z)}\right),
\qquad s(z) = 1 + \left(\frac{z}{z_R}\right)^{\!2},
\label{eq:gaussian-beam}
\end{equation}
where the on-axis intensity falls as $1/s(z)$ and the transverse profile is Gaussian with the $z$-dependent width $w(z)^2 = w_0^2\,s(z)$. The potential is $-U_0$ at the focus and rises \emph{monotonically} to $0$ far away in every direction: the well has \emph{finite depth} $U_0$ and -- important for the escape theory below -- no barrier maximum, only a plateau. The force $\mathbf F = -\nabla U$ remains conservative; Phase 4's new physics is the finite depth, not non-conservativity (the scattering push of \S\ref{sec:phase3} is a separate layer that can be recombined later).

Expanding Eq.~\eqref{eq:gaussian-beam} to second order about the focus,
\begin{equation}
U(r,z) = -U_0 + \tfrac12\!\left(\frac{4U_0}{w_0^2}\right) r^2
+ \tfrac12\!\left(\frac{2U_0}{z_R^2}\right) z^2 + \dots,
\end{equation}
identifies the harmonic stiffnesses
\begin{equation}
k_r = \frac{4U_0}{w_0^2}, \qquad k_z = \frac{2U_0}{z_R^2},
\label{eq:harmonic-limit}
\end{equation}
so the finite-depth trap \emph{reduces to the Phase-2/3 harmonic trap near the focus} -- the Phase-4 regression identity, playing the role $F_0=0$ played in Phase 3. Equation~\eqref{eq:harmonic-limit} also delivers, retroactively, the anisotropy assumed in \S\ref{sec:phase2}: $k_z/k_r = w_0^2/(2z_R^2) \ll 1$ whenever $z_R \gg w_0$, which is the paraxial regime -- the trap is soft axially because the beam spreads over the long Rayleigh range but is focused to the short waist transversely. The single dimensionless number controlling everything in this section is the well depth in thermal units,
\begin{equation}
\frac{U_0}{\kB T}.
\end{equation}

\subsection{Power parametrisation via the dipole polarizability}
\label{sec:power}

To make beam power $P$ the experimental control variable, the depth is tied to $P$ through the dipole (Rayleigh) model of the optical force \cite{JonesMaragoVolpe2015}. A sub-wavelength dielectric sphere of radius $a$ and refractive index $n_p$ in a medium of index $n_m$ has the Clausius--Mossotti polarizability
\begin{equation}
\alpha = 4\pi\varepsilon_0 n_m^2\, a^3\,\frac{m^2-1}{m^2+2},
\qquad m = \frac{n_p}{n_m},
\label{eq:polarizability}
\end{equation}
and sits in a potential $U = -\alpha I/(2c\varepsilon_0 n_m)$ proportional to the intensity $I$. With the peak intensity of a Gaussian beam of power $P$, $I_0 = 2P/(\pi w_0^2)$, and the Rayleigh range fixed by the optics, $z_R = \pi n_m w_0^2/\lambda$, the trap depth is
\begin{equation}
U_0 = \frac{\alpha I_0}{2c\varepsilon_0 n_m} = \frac{\alpha P}{\pi c\,\varepsilon_0 n_m\, w_0^2}
\;\propto\; P\,a^3\,\frac{m^2-1}{m^2+2},
\label{eq:U0-power}
\end{equation}
and the stiffnesses follow from Eq.~\eqref{eq:harmonic-limit}, so $(U_0, z_R, k_r, k_z)$ are mutually consistent by construction. Two scalings in Eq.~\eqref{eq:U0-power} define the regime: depth is \emph{linear in power}, and it falls as the \emph{cube} of the particle radius. For the reference case of this project -- a $\sim 50\,\mathrm{nm}$ nanodiamond ($n_p \approx 2.4$) in water ($n_m = 1.33$) under a $\lambda \approx 1\,\mu\mathrm{m}$ laser focused to $w_0 = 0.5\,\mu\mathrm{m}$ -- the well is only a few $\kB T$ deep at few-mW powers: exactly the marginal-trapping regime where escape and entry genuinely compete, and where the harmonic (infinite-well) idealisation of the earlier phases fails qualitatively.

\subsection{Escape: exact mean-first-passage-time theory}
\label{sec:mfpt}

Retention is quantified by the mean first-passage time (MFPT) out of the well. The escape problem is posed on the axial channel -- the on-axis cut $U(z) \equiv U(0,z) = -U_0/s(z)$, with $z$ a Cartesian coordinate -- from the focus to an absorbing point $z_b$ placed on the plateau, several $z_R$ out, where $U \approx 0$ and the particle has unambiguously left the trap.

For overdamped 1D motion $\dd z = (F/\gamma)\dd t + \sqrt{2D}\,\dd W$ with $F = -U'$, the MFPT $T(z)$ from starting point $z$ satisfies the \emph{backward} Kolmogorov equation with unit source \cite{Gardiner2009,Hanggi1990},
\begin{equation}
D\,T'' - \frac{U'}{\gamma}\,T' = -1,
\qquad T(z_b) = 0 \ \ \text{(absorbing)}, \qquad T'(z_a) = 0 \ \ \text{(reflecting)},
\label{eq:backward}
\end{equation}
where the generator acting on $T$ is the adjoint of the Fokker--Planck operator: intuitively, $T(z)$ averages the future, so it is propagated by the backward equation, and the $-1$ source accounts for the unit of time that elapses at every instant. Using $\gamma D = \kB T$ (Einstein), Eq.~\eqref{eq:backward} takes the self-adjoint (integrating-factor) form
\begin{equation}
\frac{\dd}{\dd z}\!\left[e^{-\beta U(z)}\,\frac{\dd T}{\dd z}\right] = -\frac{e^{-\beta U(z)}}{D},
\qquad \beta \equiv \frac{1}{\kB T},
\label{eq:self-adjoint}
\end{equation}
which integrates by elementary quadrature. Integrating once from the reflecting point $z_a$ (where $T'=0$) to $y$, then once more from the start point $z_0$ to the absorbing point $z_b$ (where $T=0$), gives the \emph{exact closed form}
\begin{equation}
\boxed{\;
T(z_0) = \frac{1}{D}\int_{z_0}^{z_b} e^{+\beta U(y)}
\left[\int_{z_a}^{y} e^{-\beta U(z)}\,\dd z\right] \dd y.
\;}
\label{eq:mfpt}
\end{equation}
Three properties of Eq.~\eqref{eq:mfpt} matter here. (i) It is invariant to a constant shift of $U$ -- the shift cancels between the two exponentials -- so the potential can be referenced to whatever level is numerically convenient. (ii) It requires no barrier \emph{maximum}: unlike the textbook Kramers geometry (well--saddle--well), it applies directly to the Gaussian well, which rises monotonically to a plateau. (iii) It is exact at every depth, not merely asymptotically, so it serves as ground truth for both numerical methods across the entire power sweep -- the analytic third leg of the validation triangle, restored after Phase 3 lost it.

\subsection{The deep-well limit: Kramers/Arrhenius scaling}
\label{sec:kramers}

The physical content of Eq.~\eqref{eq:mfpt} emerges in the deep-well limit $\beta U_0 \gg 1$, where the two integrands separate cleanly by scale. The inner integrand $e^{-\beta U}$ is peaked at the well bottom (where $U \approx -U_0$) and negligible on the plateau; the outer integrand $e^{+\beta U}$ is the reverse -- exponentially small in the well, of order one on the plateau. The inner integral therefore \emph{saturates} to a constant well before the outer integrand switches on, and the double integral factorises:
\begin{equation}
T \;\simeq\; \frac{Q\,R}{D},
\qquad
Q = \int e^{-\beta U(z)}\,\dd z \ \ \text{(well)},
\qquad
R = \int e^{+\beta U(y)}\,\dd y \ \ \text{(plateau)}.
\label{eq:factorised}
\end{equation}
$Q$ is the well's ``partition length'' and carries all the exponential sensitivity: with the bottom harmonic, $U \approx -U_0 + \tfrac12 k_z z^2$, the half-line integral is Gaussian,
\begin{equation}
Q = e^{\beta U_0}\int_0^\infty e^{-\beta k_z z^2/2}\,\dd z
= e^{\beta U_0}\,\sqrt{\frac{\pi \kB T}{2k_z}},
\end{equation}
while $R \approx z_b$ up to order-one corrections ($e^{\beta U}\le 1$ with the plateau at $U=0$). The result is the Kramers/Arrhenius law \cite{Kramers1940,Hanggi1990}
\begin{equation}
T \;\simeq\; \sqrt{\frac{\pi \kB T}{2k_z}}\;\frac{R}{D}\;e^{U_0/\kB T},
\label{eq:arrhenius}
\end{equation}
mean retention time exponential in the well depth, hence -- by Eq.~\eqref{eq:U0-power} -- \emph{exponential in the beam power}. The prefactor depends on where the absorbing point is placed (through $R$), but only linearly and power-independently; the $e^{U_0/\kB T}$ factor, which spans many orders of magnitude across a modest power sweep, does not. This exponential is one arm of the Phase-4 trade-off. The ratio of the exact quadrature \eqref{eq:mfpt} to the asymptote \eqref{eq:arrhenius} tends to $1$ as $U_0/\kB T \to \infty$, which is verified directly.

\subsection{Entry: the Debye--Smoluchowski capture rate}
\label{sec:capture}

The contamination channel is diffusion of a second particle from the bulk into the trap. The classical model \cite{Smoluchowski1917,Debye1942}: a reservoir at number density $n_{\rm bulk}$ far away feeds a steady diffusive current onto an absorbing sphere of radius $R_c$ (the trap core -- a particle reaching it has been captured), through the trap's own attractive radial potential $U(r) \equiv U(r,0) = -U_0\,e^{-2r^2/w_0^2}$. In spherical symmetry the steady, source-free Smoluchowski equation says the total inward current through every sphere is the same constant $C$. Writing the flux with the same integrating factor as Eq.~\eqref{eq:self-adjoint},
\begin{equation}
J(r) = -D\left(\pdv{\rho}{r} + \beta U'\rho\right)
= -D\,e^{-\beta U}\,\frac{\dd}{\dd r}\!\left(e^{\beta U}\rho\right),
\qquad
C = -4\pi r^2 J(r) = \text{const},
\end{equation}
the equation $\dd(e^{\beta U}\rho)/\dd r = C\,e^{\beta U}/(4\pi D r^2)$ integrates immediately. Applying the boundary conditions $\rho(R_c) = 0$ (absorbing core) and $\rho \to n_{\rm bulk}$, $U\to0$ as $r\to\infty$ (reservoir) fixes $C$, and the capture rate constant $k = C/n_{\rm bulk}$ is the \emph{Debye--Smoluchowski rate}
\begin{equation}
\boxed{\;
k = \frac{4\pi D}{\displaystyle\int_{R_c}^{\infty} \frac{e^{\beta U(r)}}{r^2}\,\dd r}
\quad [\mathrm{m^3\,s^{-1}}],
\;}
\label{eq:debye}
\end{equation}
which for $U=0$ reduces to the bare Smoluchowski result $k = 4\pi D R_c$ (the integral is $1/R_c$): pure diffusion onto an absorbing sphere. The structure of Eq.~\eqref{eq:debye} is a resistance-in-series formula -- each shell $[r, r+\dd r]$ contributes a diffusive resistance $e^{\beta U}/(4\pi D r^2)\,\dd r$ -- and an attractive well ($U<0$) \emph{lowers} the resistance of the shells it occupies, speeding capture. Deepening the trap therefore shortens the mean time to contamination,
\begin{equation}
T_{\rm cont} = \frac{1}{k\,n_{\rm bulk}},
\label{eq:tcont}
\end{equation}
though only algebraically -- the exponential enhancement $e^{\beta U_0}$ appears inside an integral dominated by the long free-diffusion approach, so $k$ grows far more gently with $P$ than $T_{\rm esc}$ does. That asymmetry is what makes the trade-off well-posed.

\subsection{The retention--contamination trade-off and the compromise power}
\label{sec:tradeoff}

The two closed forms \eqref{eq:arrhenius} and \eqref{eq:debye}--\eqref{eq:tcont} pull in opposite directions as functions of beam power:
\begin{itemize}[leftmargin=1.6em]
\item \textbf{Retention} $T_{\rm esc}(P) \sim A(P)\, e^{cP}$ with $c = U_0/(P\,\kB T)$ fixed by Eq.~\eqref{eq:U0-power}: rises \emph{exponentially} with power.
\item \textbf{Time-to-contamination} $T_{\rm cont}(P) = 1/[k(P)\,n_{\rm bulk}]$: falls monotonically with power, algebraically.
\end{itemize}
The useful single-particle measurement window ends when \emph{either} the held particle escapes \emph{or} a second one arrives, i.e.\ after a time of order
\begin{equation}
T_{\rm hold}(P) = \min\bigl(T_{\rm esc}(P),\, T_{\rm cont}(P)\bigr).
\end{equation}
Since one branch increases and the other decreases, the minimum is maximised exactly where the two curves cross, defining the \emph{compromise power} $P^{*}$:
\begin{equation}
T_{\rm esc}(P^{*}) = T_{\rm cont}(P^{*}).
\label{eq:pstar}
\end{equation}
Below $P^{*}$ the trap is under-powered and the particle is lost to Brownian escape before contamination matters; above it the trap is over-powered and a contaminant arrives first. Because $T_{\rm esc}$ crosses steeply (exponentially) while $T_{\rm cont}$ moves gently, $P^{*}$ is sharply defined, and its dependence on the bulk concentration is only logarithmic: from Eq.~\eqref{eq:pstar}, $c P^{*} \approx -\ln n_{\rm bulk} + \text{const}$, so even order-of-magnitude changes in contamination level shift the optimal power modestly. For the reference nanodiamond suspension ($n_{\rm bulk} = 2\times10^{13}\,\mathrm{m^{-3}}$), the crossover sits at $P^{*} \approx 11\,\mathrm{mW}$ (Figure~\ref{fig:phase4}).

\begin{figure}[h]
\centering
\includegraphics[width=0.9\textwidth]{figures/phase4_validation.png}
\caption{The Phase-4 trade-off (\texttt{scripts/validate\_phase4.py}): retention time $T_{\rm esc}(P)$ (rising exponentially, Eq.~\eqref{eq:arrhenius}) against time-to-contamination $T_{\rm cont}(P)$ (falling, Eqs.~\eqref{eq:debye}--\eqref{eq:tcont}) across the beam-power sweep for a $\sim50\,\mathrm{nm}$ nanodiamond in water. The single-particle holding window $\min(T_{\rm esc}, T_{\rm cont})$ is maximised at the crossover, the compromise power $P^{*}$.}
\label{fig:phase4}
\end{figure}

\subsection{Numerical treatment in brief}

Each closed form above is shadowed by one or two independent numerical routes, restoring the full three-way validation. \emph{Escape}: (i) the backward equation \eqref{eq:self-adjoint} is discretised directly in its self-adjoint form -- face weights $e^{-\beta U}$, the exact continuum analogue of Chang--Cooper's exponential fitting, hence stable at any barrier steepness -- and solved as a single symmetric tridiagonal system, with no time-stepping; (ii) as a cross-check, the forward equation is time-stepped from a source at the well bottom with an \emph{absorbing} outer wall (the one boundary-condition change the entry/exit problem forces on the Phase-1 solver), and the retention time recovered from the survival probability as $\avg{t_{\rm esc}} = \int_0^\infty S(t)\,\dd t$, $S(t)=\int\rho\,\dd z$; (iii) Brownian dynamics gains an absorbing test and first-passage-time accounting per particle. Direct BD sampling of a deep well is an exponentially rare-event problem, so the BD leg is run at moderate depth, where all three routes overlap; conversely the tridiagonal backward solve becomes ill-conditioned past $U_0/\kB T \sim 20$ (its weights span the $e^{\beta U_0}$ dynamic range), where the unconditionally stable analytic quadrature \eqref{eq:mfpt} carries the deep end -- the methods hand over to each other across the sweep. \emph{Entry}: the steady capture problem, written as $\dd/\dd r\,[\,r^2 e^{-\beta U}\,\dd(e^{\beta U}\rho)/\dd r\,] = 0$, is again a self-adjoint tridiagonal solve, with the rate read off the converged inner-face current and cross-checked against Eq.~\eqref{eq:debye}.

% ======================================================================
\section{Validation summary}
\label{sec:validation}
% ======================================================================

The numerical evidence is summarised once, compactly; the scripts \texttt{validate\_phase$N$.py} ($N=1,\dots,4$, in \texttt{scripts/}) regenerate every number and figure, and the test suite (\texttt{tests/}) enforces them continuously.

\begin{itemize}[leftmargin=1.6em]
\item \textbf{Phase 1.} Both methods reproduce the analytic OU relaxation and Boltzmann stationary state: stationary $\sigma_x$ to $0.28\%$ (Brownian dynamics) and to numerical precision (Fokker--Planck); probability conserved to $\sim10^{-12}$; density non-negative throughout.
\item \textbf{Phase 2.} Both methods reproduce the analytic moments $\avg{r^2}=2\sigma_r^2$ and $\avg{z^2}=\sigma_z^2$ and the Rayleigh and Gaussian marginals of \S\ref{sec:rayleigh} to within the $\sim3\%$ cross-check tolerance (FP much tighter); the stationary density conserves probability ($\sim10^{-13}$), stays non-negative, and factorises as the separable potential requires.
\item \textbf{Phase 3.} $F_0=0$ reproduces Phase 2 exactly, bit-for-bit in Brownian dynamics and to $10^{-6}$ in the Fokker--Planck solve; the Rayleigh invariant of \S\ref{sec:rayleigh-invariant} holds for $F_0>0$; the steady current is nonzero and circulating, with a normalised discrete divergence of $\sim2\times10^{-4}$ against a tolerance of $5\times10^{-2}$; BD and FP agree on moments and marginals to $\sim0.3\%$.
\item \textbf{Phase 4.} The finite-depth potential's curvature reproduces $k_r, k_z$ (harmonic-limit regression, Eq.~\eqref{eq:harmonic-limit}); the FP backward-equation MFPT matches the exact quadrature \eqref{eq:mfpt} to better than $1\%$ across the full power sweep, and the survival-decay route and BD first-passage statistics agree at moderate depth (BD within Monte-Carlo error); the FP capture rate matches Debye--Smoluchowski \eqref{eq:debye} to better than $0.5\%$, and reduces to $4\pi D R_c$ at $U=0$; both escape routes approach the Arrhenius asymptote \eqref{eq:arrhenius} as $U_0/\kB T$ grows.
\end{itemize}

% ======================================================================
\section{Discussion}
% ======================================================================

The arc of the four phases is a controlled retreat from exact solvability, with the validation structure engineered so that trust is never asked for without collateral. Phases 1--2 exercise every numerical component -- exponential-fitting fluxes, unconditionally stable implicit time-stepping, the cylindrical measure and its staggered-grid regularity, dimensional splitting -- in settings where a closed-form answer grades them. Phase 3 then poses a question with no closed-form answer at all: the non-conservative steady state, whose defining observable is a circulating probability current rather than a density formula. There the two-method agreement, the exact Rayleigh invariant of \S\ref{sec:rayleigh-invariant}, and the discrete divergence condition together substitute for the missing analytic leg. Phase 4 shows that pushing the model \emph{toward} realism can also restore solvability: making the well finite-depth -- unavoidable if escape and entry are to exist at all -- brings with it the exact first-passage quadrature \eqref{eq:mfpt} and the Debye--Smoluchowski rate \eqref{eq:debye}, so the applied headline result (the compromise power $P^{*}$) rests on a three-way check at both of its arms.

The closed-form structure of \S\ref{sec:phase4} is also what makes the applied result robust. The exponential $e^{U_0/\kB T}$ in Eq.~\eqref{eq:arrhenius} dominates every prefactor uncertainty (including the placement of the absorbing boundary, which enters only through the plateau length $R$), and the logarithmic sensitivity of $P^{*}$ to $n_{\rm bulk}$ means the compromise power is a property of the particle and optics far more than of the (poorly known) contamination level. The natural extensions are structural rather than conceptual: the full two-dimensional axisymmetric first-passage problem in place of the 1D axial and spherical radial reductions; recombining the Phase-3 scattering push with the finite-depth well, where escape becomes a non-equilibrium first-passage problem and the analytic leg is lost again -- but now with two independently validated numerical methods ready for it; and replacing the dipole-model force field with T-matrix-computed optical forces.

% ======================================================================
\section{Conclusion}
% ======================================================================

Starting from the inertial Langevin equation for a bead in a viscous fluid, this report derived the overdamped Ornstein--Uhlenbeck description of a particle in a harmonic optical trap and its exact solution, together with the equivalent Fokker--Planck description and its Boltzmann stationary state (Phase 1). The theory was extended to the axisymmetric trap geometry, where Cartesian separability reduces the dynamics to three decoupled OU processes while the cylindrical measure turns the radial marginal into an exact Rayleigh distribution (Phase 2). Adding the scattering force produced a force field with nonzero curl, breaking detailed balance: the stationary state becomes a non-equilibrium steady state carrying a circulating, divergence-free probability current -- the Brownian vortex -- with no closed-form density, and the radial Rayleigh marginal surviving as the one exact invariant (Phase 3). Replacing the harmonic well with the finite-depth Gaussian-beam potential opened the trap to escape and entry, both governed by closed-form first-passage theory: the exact mean-first-passage-time quadrature with its deep-well Kramers/Arrhenius limit $T_{\rm esc} \propto e^{U_0/\kB T}$, and the Debye--Smoluchowski capture rate (Phase 4). Their opposing dependence on beam power yields the project's applied result: the single-particle holding window is maximised at the compromise power where retention time and time-to-contamination cross. At every stage the analysis was carried by two independent numerical methods validated against each other and against every closed form available -- a discipline that pays for itself exactly where the closed forms run out.

% ======================================================================
\begin{thebibliography}{9}

\bibitem{JonesMaragoVolpe2015}
P.~H.~Jones, O.~M.~Mara\`{g}\`{o}, G.~Volpe,
\emph{Optical Tweezers: Principles and Applications},
Cambridge University Press (2015).

\bibitem{VolpeVolpe2013}
G.~Volpe, G.~Volpe,
\emph{Simulation of a Brownian particle in an optical trap},
Am.\ J.\ Phys.\ \textbf{81}, 224 (2013).

\bibitem{ChangCooper1970}
J.~S.~Chang, G.~Cooper,
\emph{A practical difference scheme for Fokker--Planck equations},
J.\ Comput.\ Phys.\ \textbf{6}, 1 (1970).

\bibitem{Gardiner2009}
C.~Gardiner,
\emph{Handbook of Stochastic Methods for Physics, Chemistry and the Natural Sciences},
4th ed., Springer (2009).

\bibitem{PeacemanRachford1955}
D.~W.~Peaceman, H.~H.~Rachford,
\emph{The numerical solution of parabolic and elliptic differential equations},
J.\ Soc.\ Ind.\ Appl.\ Math.\ \textbf{3}, 28 (1955).

\bibitem{Roichman2008}
Y.~Roichman, B.~Sun, A.~Yevick, J.~Amato-Grill, D.~G.~Grier,
\emph{Influence of nonconservative optical forces on the dynamics of optically trapped colloidal spheres: the fountain of probability},
Phys.\ Rev.\ Lett.\ \textbf{101}, 128301 (2008).

\bibitem{Sun2009}
B.~Sun, J.~Lin, E.~Darby, A.~Y.~Grosberg, D.~G.~Grier,
\emph{Brownian vortexes},
Phys.\ Rev.\ E \textbf{80}, 010401(R) (2009).

\bibitem{Kramers1940}
H.~A.~Kramers,
\emph{Brownian motion in a field of force and the diffusion model of chemical reactions},
Physica \textbf{7}, 284 (1940).

\bibitem{Hanggi1990}
P.~H\"anggi, P.~Talkner, M.~Borkovec,
\emph{Reaction-rate theory: fifty years after Kramers},
Rev.\ Mod.\ Phys.\ \textbf{62}, 251 (1990).

\bibitem{Smoluchowski1917}
M.~von~Smoluchowski,
\emph{Versuch einer mathematischen Theorie der Koagulationskinetik kolloider L\"osungen},
Z.\ Phys.\ Chem.\ \textbf{92}, 129 (1917).

\bibitem{Debye1942}
P.~Debye,
\emph{Reaction rates in ionic solutions},
Trans.\ Electrochem.\ Soc.\ \textbf{82}, 265 (1942).

\end{thebibliography}

\appendix
% ======================================================================
\section{Symbol reference}
% ======================================================================

\begin{table}[h]
\centering
\begin{tabular}{@{}lll@{}}
\toprule
Symbol & Meaning & Units \\
\midrule
$a$ & particle radius & m \\
$\eta$ & fluid dynamic viscosity & Pa\,s \\
$T$ & temperature & K \\
$\gamma$ & Stokes drag coefficient, $6\pi\eta a$ & kg/s \\
$D$ & diffusion coefficient, $\kB T/\gamma$ & m$^2$/s \\
$\beta$ & inverse thermal energy, $1/\kB T$ & J$^{-1}$ \\
$k$, $k_r$, $k_z$ & trap (harmonic) stiffnesses & N/m \\
$\tau$, $\tau_r$, $\tau_z$ & relaxation times, $\gamma/k$ & s \\
$\sigma_x$, $\sigma_r$, $\sigma_z$ & stationary standard deviations, $\sqrt{\kB T/k}$ & m \\
$U(\mathbf x)$ & trap potential & J \\
$P$, $\rho$ & probability density & m$^{-1}$, m$^{-3}$ \\
$\mathbf J$ & probability current & (density)\,m\,s$^{-1}$ \\
$F_0$, $w_0$ & scattering-push amplitude and transverse width (Phase 3) & N, m \\
$w_0$, $z_R$ & beam waist and Rayleigh range (Phase 4) & m \\
$U_0$ & finite well depth; $U_0/\kB T$ the depth in thermal units & J \\
$\alpha$ & dipole (Clausius--Mossotti) polarizability & C\,m$^2$/V \\
$P$ (in \S\ref{sec:phase4}) & beam power & W \\
$T(z)$ & mean first-passage time from $z$ & s \\
$k$ (in \S\ref{sec:capture}) & diffusion-limited capture rate constant & m$^3$/s \\
$n_{\rm bulk}$ & bulk number density of contaminant particles & m$^{-3}$ \\
$P^{*}$ & compromise beam power, $T_{\rm esc}=T_{\rm cont}$ & W \\
$\mathrm{Pe}$ & local cell P\'eclet number, $v\Delta x/D$ & --- \\
$B(z)$ & Bernoulli function $z/(e^z-1)$ (Chang--Cooper weights) & --- \\
\bottomrule
\end{tabular}
\caption{Symbols used throughout this report, matching the conventions of \texttt{src/params.py}. Note the reuse of $k$ and $P$ across contexts, disambiguated in the text.}
\end{table}

\end{document}
